% Part:sets-functions-relations
% Chapter: size-of-sets
% Section: introduction

\documentclass[../../../include/open-logic-section]{subfiles}

\begin{document}

\olfileid{sfr}{siz}{int}

\olsection{引言}

19世纪70年代，当 Georg Cantor（康托尔）发展集合论时，他的目标之一是使无限集合的观念变得易于接受——即中世纪哲学家所说的“实际无限”。这其中的一个关键部分，是他对不同集合\emph{大小}的处理。如果 $a$、$b$ 和 $c$ 互不相同，那么集合 $\{a, b, c\}$ 直观上就比 $\{a, b\}$ \emph{更大}。但无限集呢？它们的大小都一样吗？事实并非如此。

这里的首要重要概念是枚举。我们可以通过列出其所有元素来枚举任何有限集。对于某些无限集，只要允许列表本身是无限的，我们也能列出其所有元素。这样的集合被称为\emph{可数的}。Cantor 的一个惊人结果（我们到本章末尾就会完全理解）是：有些无限集是不可数的。

\end{document}